\begin{abstract}
  W pracy zbadano wpływ topologii systemów agentowych na jakość odpowiedzi
  na pytania wymagające analizy dokumentów ubezpieczeniowych -- Ogólnych Warunków
  Ubezpieczenia (OWU) oraz dokumentów zawierających informacje o~produkcie
  ubezpieczeniowym (IPID).  Środowisko testowe wybrano ze względu na jego złożoność: specjalistyczny
  język prawniczo-ubezpieczeniowy łączy się z~wielopoziomową strukturą warunków
  i~wyłączeń rozproszonych w~wielu paragrafach. Wymagało to od badanych
  architektur skutecznego wykorzystania mechanizmów wyszukiwania i~wnioskowania.

  Porównano sześć architektur opartych na dużych modelach językowych
  (\textit{ang.}~\textit{Large Language Models}, LLM): deterministyczną,
  planer--wykonawca, router--specjalista, tablicową, hierarchiczną oraz ReAct.
  Bazę wiedzy systemu stanowiły 22 polskie dokumenty ubezpieczeniowe pochodzące
  od trzech głównych towarzystw ubezpieczeniowych. Opracowano autorski zestaw 90
  pytań w~czterech poziomach trudności: bardzo łatwym, łatwym, trudnym i~bardzo
  trudnym -- każde pytanie wymagało odniesienia się do treści zgromadzonych
  dokumentów. Łącznie przeprowadzono 540 uruchomień, które oceniono według
  wspólnego protokołu ewaluacyjnego uwzględniającego wierność, poprawność,
  zwięzłość, trafność pobieranych dokumentów, opóźnienie oraz zużycie tokenów.

  Wyniki wskazują, że większa złożoność orkiestracji nie gwarantuje wyższej
  jakości odpowiedzi. Architektura router--specjalista osiąga najwyższą średnią
  wierność (0{,}811), a~deterministyczna -- najwyższą średnią poprawność
  (0{,}716), obie przy koszcie porównywalnym z~najprostszymi układami.
  Architektury hierarchiczna i~ReAct zużywają kilkukrotnie więcej tokenów bez
  proporcjonalnego przyrostu jakości -- analiza jakościowa wskazuje, że dłuższe
  ścieżki przetwarzania prowadzą do utraty kontroli nad kontekstem i~propagacji
  błędu między modułami. Sformułowano wnioski praktyczne dotyczące doboru
  architektury do przypadku użycia, opłacalności złożoności orkiestracji oraz
  roli przygotowania danych jako czynnika warunkującego skuteczność systemu
  niezależnie od topologii.

  \bigskip
  \noindent\textbf{Słowa kluczowe:} architektury agentowe, duże modele językowe, ewaluacja systemów, generowanie wspomagane wyszukiwaniem, systemy wieloagentowe
\end{abstract}
